% !PRINTED 172
\weberpage{172}
%<<CONT
Substitution, so sind diese Formen homogene Functionen von
$n$ Variablen, und es können also nicht mehr als $n$ von einander
unabhängige existiren. Damit ist nicht gesagt, dass sich alle
diese Formen \gesperrt{rational} durch $n$ unter ihnen ausdrücken lassen,
aber zwischen $n+1$ Invarianten muss immer eine rationale
Gleichung bestehen, die sich durch Elimination der $n$ Variablen
ergiebt.

Es ist aber noch die umgekehrte Frage zu untersuchen, ob
es wirklich für eine lineare Substitutionsgruppe von der Dimension
$n$ immer $n$ unabhängige Invarianten giebt.

Wenn wir für die Variablen feste Werthe setzen, so erhalten
dadurch die sämmtlichen Invarianten der Gruppe gleichfalls
bestimmte Werthe. Die Frage, die wir noch zu beantworten
haben, ist nun die, ob zu einem Werthsysteme der Invarianten
bestimmte Werthsysteme der Variablen, und zwar in endlicher
Anzahl, existiren. Wenn dies bewiesen ist, so können wir die
Variablen als (mehrwerthige) algebraische Functionen der Invarianten
auffassen, und die Anzahl der von einander unabhängigen
Invarianten kann nicht kleiner sein, als die Anzahl
der Variablen, weil sonst ein Theil der Variablen, wenn die
Werthe der Invarianten gegeben sind, noch willkürlich bleiben
würde.

Wenn $(x)$ ein einem bestimmten Werthsysteme der Invarianten
entsprechendes Werthsystem der Variablen ist, und $A$ eine Substitution
der Gruppe $S$, so entspricht nach der Natur der Invarianten
das System $A(x)$ demselben Werthsysteme der Invarianten,
und unser Problem hat also mindestens so viele
Lösungen, als der Grad der Gruppe beträgt. Dass dies aber die
genaue Anzahl der Lösungen ist, geht aus den folgenden Betrachtungen
hervor.

In besonderen Fällen, d. h. für besondere Werthe der Invarianten,
können von diesen Werthsystemen der Variablen
mehrere zusammenfallen. Dies kann aber nur für solche Werthsysteme
der $(x)$ geschehen, für die eine Relation von der Form
$(x) = A(x)$ besteht, denn aus $A(x) = B(x)$ würde $(x) = A^{-1}B(x)$
folgen, was in der Form $(x) = A(x)$ enthalten ist.

Wir nehmen eine lineare homogene Function $\Theta$ der Variablen
$x_1, x_2, \dots x_n$ an, die wir so bezeichnen:
\weq{1}{\Theta = \Theta(x) = c_1 x_1 + c_2 x_2 + \cdots + c_n x_n.}
